\documentclass[12pt,letterpaper]{article}
\usepackage[utf8]{inputenc}
\usepackage[spanish]{babel}
\usepackage{amsmath}
\usepackage{amsfonts}
\usepackage{amssymb}
\usepackage[margin=2.5cm]{geometry}

\begin{document}

\begin{center}
    \Large \textbf{Evaluación de Examen 2: Unidad II} \\
    \large Física para Ciencias de la Salud (005-1713) \\
    \vspace{0.5cm}
    \textbf{Estudiante:} Estefani Santoya \quad \textbf{C.I.:} 31.686.375
    \vspace{0.3cm} \\
    \large \textbf{Calificación Final: 10/10 puntos}
\end{center}

\vspace{0.5cm}

\section*{Análisis de Resultados}

\subsection*{1. Cinemática (Aceleración media)}
\textbf{Procedimiento y resultado:} Plantea correctamente la fórmula y realiza la sustitución de las velocidades y el tiempo ($0,6 \text{ m/s} / 0,15 \text{ s}$). El resultado matemático es exacto ($4 \text{ m/s}^2$) y las unidades del Sistema Internacional están debidamente expresadas. \\
\textbf{Veredicto:} \textbf{Correcto} (2/2 pts).

\subsection*{2. Dinámica (Leyes de Newton)}
\textbf{Procedimiento y resultado:} Extrae con orden la masa y la fuerza, para luego hacer el despeje algebraico perfecto partiendo de la Segunda Ley de Newton ($a = \frac{F}{m}$). Sustituye los valores dividiendo $150 \text{ N}$ entre $25 \text{ kg}$ y halla el resultado correcto de $6 \text{ m/s}^2$. \\
\textbf{Veredicto:} \textbf{Correcto} (2/2 pts).

\subsection*{3. Trabajo Mecánico}
\textbf{Procedimiento y resultado:} El análisis teórico es sobresaliente. En su sección de datos, la estudiante justifica que tanto la fuerza como el desplazamiento tienen la misma dirección ascendente, por lo que aplica un ángulo de $0^\circ$ e indica explícitamente que $\cos(0^\circ) = 1$. La multiplicación se ejecuta sin errores para obtener el valor de $320 \text{ Joules}$. \\
\textbf{Veredicto:} \textbf{Correcto} (2/2 pts).

\subsection*{4. Energía Potencial}
\textbf{Procedimiento y resultado:} Extrae los datos a la perfección y plantea la ecuación de energía potencial gravitatoria ($E_p = m \cdot g \cdot h$). El producto secuencial ($1,2 \text{ kg} \cdot 9,8 \text{ m/s}^2 \cdot 1,5 \text{ m}$) brinda la cifra matemáticamente exacta y comprobada de $17,64 \text{ Joules}$. \\
\textbf{Veredicto:} \textbf{Correcto} (2/2 pts).

\subsection*{5. Potencia Mecánica}
\textbf{Procedimiento y resultado:} Reconoce y emplea directamente la fórmula de la potencia mecánica como la relación del trabajo entre el tiempo ($P = \frac{W}{\Delta t}$). Ejecuta la sustitución de $75 \text{ J} / 60 \text{ s}$ logrando el resultado perfecto de $1,25 \text{ Watts}$. \\
\textbf{Veredicto:} \textbf{Correcto} (2/2 pts).

\vspace{0.5cm}
\section*{Comentarios Generales y Sugerencias}
¡Muchas felicidades por un examen perfecto! El nivel de desempeño mostrado por la estudiante es digno de reconocimiento.
\begin{itemize}
    \item La subdivisión de los problemas en ``Datos'', ``Fórmula'', ``Solución'' y ``Resultado'' (escrito de forma analítica en una oración) es la estructura metodológica ideal en el área científica y clínica.
    \item Es loable el dominio mostrado sobre las unidades del Sistema Internacional, integrando de manera coherente las magnitudes en todas las fases del desarrollo de las ecuaciones (como usar correctamente $\text{m/s}^2$ en la aceleración o transformar de $N \cdot m$ a Joules).
    \item Se invita a la estudiante a mantener este estupendo estándar de excelencia en sus próximas evaluaciones.
\end{itemize}

\end{document}